% ============================================================
% SECTION 2: BACKGROUND AND RELATED WORK. Written 2026-07-27.
%
% Job, per PAPER_STRUCTURE.md and the HANDOFF: introduce only what the rest of
% the paper uses; define "removable redundancy", which Section 1 uses as a
% technical term and which nothing else in the paper defines; and hold the
% contemporary instances of the missing control, with the arithmetic (decision
% D1). The literature is motivation and related work here, never the spine.
%
% Mohammad, 2026-07-27: Socievole and Pizzuti are cited as current practice,
% briefly, with no tables and no numbers. This is a related-work section.
%
% Every characterization below comes from a paper read in full, and the reading
% is recorded in PAPER_RESTRUCTURE/references/:
%   NOTES_satuluri2011.md          the 4.3 transfer control and the fixed-k arms
%   NOTES_chen_and_fastcd.md       Chen (what they own, what they do not) and
%                                  Laeuchli (the regime where sparsification pays)
%   NOTES_effres_2026.md           Pari et al., football 0.87 against max 0.6046
%   NOTES_kcore_ijain2026.md       Setiadi et al., karate 0.63 against max 0.4198
%   NOTES_socievole_pizzuti_2026.md  read in full; cited here as practice only
%   P1_EVIDENCE.md                 Hamann, Blagus, sampling line, IJCAI survey
%   P2_BACKGROUND_RESEARCH.md      metric backbone, Maslov-Sneppen, Graclus/MLR-MCL
% STYLE: no em-dashes, no AI-characteristic phrasing (see STYLE.md).
% ============================================================

\subsection{Detection algorithms and their degrees of freedom}

Community detection algorithms divide into two families, and the division is the
first of the two conditions this paper identifies.

One family chooses its own number of communities. Modularity
optimizers~\cite{newman2004finding} such as Louvain~\cite{blondel2008fast} and
Leiden~\cite{traag2019louvain} maximize \eqref{eq:modularity} by agglomeration,
stopping when no move improves the objective. Infomap~\cite{rosvall2008maps}
minimizes the description length of a random walk, and label
propagation~\cite{raghavan2007near} iterates a local majority rule to a fixed
point. None takes the number of communities as an input, and each returns
whatever count its objective and its search prefer. Modularity in particular has
a resolution limit~\cite{fortunato2007resolution}: at $\gamma = 1$ it cannot
resolve communities below a size set by the total number of edges, so the count
it returns is a property of the graph's size as much as of its structure.

The other family takes the number of communities as an input. Multilevel graph
partitioners such as Metis~\cite{karypis1998metis} and
Graclus~\cite{dhillon2007weighted} coarsen the graph, partition the coarsest
level into a requested $k$ parts, and refine on the way back up, subject to a
balance constraint on part sizes. Flow-based clustering with a coarsening
schedule, as in MLR-MCL~\cite{satuluri2009scalable}, sits between the two: the
count follows from an inflation parameter rather than from a requested $k$.

The difference matters here because sparsification changes what a
free-granularity algorithm returns, and does not change what a fixed-$k$
algorithm returns. Section~\ref{sec:protocol} makes that a control and
Section~\ref{sec:experiments} makes it a result.

\subsection{Sparsification and what it is guaranteed to preserve}

Spectral sparsification is the part of this literature with proofs. Benczúr and
Karger~\cite{benczur2015randomized} preserve all cuts to within $1 \pm \epsilon$
by sampling edges according to connectivity, and Spielman and
Srivastava~\cite{spielman2008graph} preserve the Laplacian quadratic form by
sampling according to effective resistance. The guarantee is about the Laplacian,
not about communities, and Section~\ref{sec:intro} states why the two are related
without being the same.

The methods actually used before community detection are cheaper and carry no
such guarantee. Similarity-based methods rank each edge by the overlap of its
endpoints' neighbourhoods, as in the local sparsification of Satuluri et
al.~\cite{satuluri2011local}. Degree-based methods score an edge from its
endpoints' degrees alone, as in DSpar~\cite{liu2023dspar}, which approximates
effective resistance by $1/d_u + 1/d_v$. Backbone extraction keeps edges that are
statistically surprising against a local null, as in the disparity
filter~\cite{serrano2009extracting} and its noise-corrected
successor~\cite{coscia2017backboning}. Learned methods train a model to choose the
subgraph~\cite{wu2020gsgan}. Hamann et al.~\cite{hamann2016structure} compared
these families systematically on social networks and released the
implementations that are now distributed in standard graph libraries.

Two recent surveys bound the field. Chen et al.~\cite{chen2024demystifying}
benchmark twelve sparsifiers against sixteen graph properties on fourteen
networks and conclude that no single sparsifier preserves all of them and that
the choice should follow the downstream task. They also establish, at benchmark
scale, that pruning fragments graphs and that fidelity to the full graph's
partition degrades monotonically with the prune rate. Both facts are theirs and
this paper does not claim them. Hashemi et al.~\cite{hashemi2024survey} survey
the wider reduction field of sparsification, coarsening and condensation, and
their task list is a useful negative: the modern literature has reoriented around
graph neural network prediction, and community detection has largely dropped out
of it as an evaluation target.

\subsection{Removable redundancy}\label{sec:background:redundancy}

The second condition in this paper's claim is that a graph carry edges whose
removal does not destroy its community structure. The quantity has an established
formal relative. The \emph{metric backbone} of a weighted graph is the union of
its all-pairs shortest paths, so every edge outside it is redundant in the exact
sense that no shortest path uses it. Correia et
al.~\cite{correia2023contact} report that the backbone is a small fraction of the
edges of real contact networks, and Dreveton et al.~\cite{dreveton2024metric}
prove that deleting everything outside it leaves community structure intact on a
broad class of weighted random graphs with communities.

Removable redundancy in the sense used here is broader than the metric one, since
the edges a sparsifier may delete without cost are not only those absent from
shortest paths, and we do not have a procedure that measures it directly. What
the results of Section~\ref{sec:experiments:boundary} establish is that the
quantity varies across graphs, that average degree is its strongest measured
correlate, and that degree heterogeneity is not a correlate at all.

\subsection{Sparsify, then detect}

Satuluri et al.~\cite{satuluri2011local} introduced the pipeline this paper
examines and reported speedups commonly between $10\times$ and $50\times$ with
little or no loss of cluster quality, together with improved accuracy for two of
the four clustering algorithms they tested. Those two were the partitioners that
take $k$ as an input.

The practice spread. Reference implementations entered graph
libraries~\cite{hamann2016structure}, sparsification was applied before community
detection at industrial scale~\cite{satuluri2020simclusters}, generative models
were trained to produce the subgraph~\cite{wu2020gsgan}, and triangle-aware
spectral sparsifiers were proposed specifically because sparsification appeared
to be a route to community detection~\cite{sotiropoulos2021triangle}. It remains
current: Socievole and Pizzuti weight a graph by effective resistance, apply a
kernel, sparsify, and then detect communities on what
remains~\cite{socievole2024community, socievole2026effective}, and further
effective-resistance and $k$-core variants of the same pipeline have appeared
since~\cite{pari2026effective, setiadi2025community}.

A parallel line asks whether sampling preserves community structure at all.
Leskovec and Faloutsos~\cite{leskovec2006sampling} set the sampling problem,
Maiya and Berger-Wolf~\cite{maiya2010sampling} sample explicitly for community
structure, and Blagus et al.~\cite{blagus2015sampling} report that sampling can
manufacture the appearance of community structure rather than preserve it. Peng
et al.~\cite{peng2014accelerating} reduce the graph by $k$-core before detecting,
which is the coarsening-flavoured version of the same idea and removes vertices
as well as edges.

\subsection{How the pipeline is evaluated}

The evaluation question that Section~\ref{sec:protocol} formalizes was raised in
the paper that introduced the pipeline. Satuluri et al.\ state in their \S4.3
that the clusters obtained from a sparsified graph cannot be scored on that same
graph, because doing so says nothing about how well the sparsification retained
the structure of the original, and they report every conductance on the original
graph accordingly. They also fix the number of communities once per network and
supply it to every arm, so their comparisons are matched on granularity by
construction, and they charge the sparsification time to the speedup. On all
three counts the original paper is stricter than much of what followed it.

What followed is less uniform. In the released code of the most thorough public
benchmark of sparsification, modularity for the sparsified condition is computed
on the sparsified graph~\cite{chen2024demystifying}. Two recent
sparsify-then-detect papers report objective values that exceed what their graphs
admit: an effective-resistance method reports modularity $0.87$ on the American
College Football network~\cite{girvan2002community}, whose maximum over all
partitions is $0.6046$, and a $k$-core method reports $0.63$ on Zachary's karate
club~\cite{zachary1977information}, whose maximum is
$0.4198$~\cite{pari2026effective, setiadi2025community}. A modularity above the
graph's maximum is not a matter of interpretation; it identifies the graph being
scored as the sparsified one. Section~\ref{sec:experiments:accounting} measures
what that choice is worth.

Two further hazards are documented prior art rather than findings of this paper.
Hamann et al.~\cite{hamann2016structure} warn that normalized mutual information
reports higher similarity for larger numbers of communities, which makes it
unsuitable for comparing partitions of sparsified networks, and point to
chance-corrected indices~\cite{vinh2010information} in its place. Gottesb\"{u}ren
et al.~\cite{gottesburen2025linear} observe that the cost of sparsifying can
cancel the speedup it is meant to produce. Separating a degree-driven effect from
a community-driven one requires a null that holds the degree sequence fixed, for
which the standard instrument is the degree-preserving edge
swap~\cite{maslov2002specificity} over configuration-model
graphs~\cite{molloy1995critical}.

Finally, one paper marks the boundary of this paper's scope rather than falling
inside it. Laeuchli~\cite{laeuchli2020fast} analyses community detection on
sparsified graphs for spectral solvers on dense planted-partition models, which
is the regime in which sparsification genuinely pays, because the detection cost
there is superlinear in the number of edges and the planted structure survives
heavy deletion. The graphs and detectors examined here are the ones the pipeline
is applied to in practice, and they are not that regime.
